\scanpage{5}\mspage{3}
pour \unsure{tout|tous} $a\in M_0$, donc $M_0$, muni des $\pi_0(a,b)$ et de leur composition, devient une \uline{catégorie}, \deleted{\unsure{de}} soit $M_{\pi}$ ; et les \deleted{\unsure{éléments unité de}} \inserted{flèches identiques} de celle-ci doivent jouer ce rôle d'unité bilatère également pour les ``Hom externe'' $\mathrm{Ch}_{\infty}(a,b)$, à valeurs dans $\mathrm{Hot}\unsure{. |{}_{0}}$

On peut dire que pour $a=a_0$ fixé, les \deleted{$\mathrm{Hom}$} $\deleted{\unsure{\mathrm H}}\mathrm{Ch}_{\infty}(a_0,b)$ ($b$ variable) forment les valeurs d'un \uline{foncteur} covariant
\[
M_{\pi}\unsure{\longrightarrow|\longleftrightarrow}\mathrm{Hot}\qquad b\mapsto\mathrm{Ch}_{\infty}(a_0,b)\unsure{,|;}
\]
\deleted{\unsure{Le fait}} en associant à
\[
b\xrightarrow{u}b'\qquad u\in\pi_0(b,b')
\]
la flèche
\[
\mathrm{Ch}_{\infty}(a,b)\unsure{\text{---}|\longrightarrow}\mathrm{Ch}_{\infty}(b')
\]
\unsure{correspondante} $\unsure{*_1|\alpha_1}$, \unsure{Que|que} ce soit bien là une loi fonctorielle doit résulter formellement de l'associativité des accouplements sur les $\mathrm{Ch}_{\infty}$. De $\unsure{\overline{\mathrm m}|=}$, on a une loi fonctorielle, pour $b$ fixé
\[
M_{\pi}^{o}\longrightarrow\mathrm{Hot}\qquad a\unsure{\vdash|\mapsto}\mathrm{Ch}_{\infty}(a,b),
\]
grâce à l'associativité. En fait, l'associativité impliquera aussi qu'on a un \uline{bifoncteur}
\[
\begin{aligned}
M_{\pi}^{o}\times M_{\pi}&\longrightarrow\mathrm{Hot}\\
(a,b)&\longrightarrow\deleted{\unsure{\mathrm H}}\mathrm{Ch}_{\infty}(a,b),
\end{aligned}
\]
% archive sheet number: 4

